% !TeX spellcheck = pt_BR
\documentclass[12pt, t]{beamer}
% Para alterar a linguagem do documento
\usepackage[portuges]{babel}
% Para aceitar caracteres especias deretamente do teclado
\usepackage[utf8]{inputenc}
% Para incluir figuras
\usepackage{graphicx}
% Para melhor ajuste da posisao das figuras
\usepackage{float}
% Para ajustar as dimensoes do layout da pagina
\usepackage{geometry}
% Para formatar estrutura e informacoes de formulas matematicas
\usepackage{amsmath}
% Para incluir simbolos especiais em formulas matematicas
\usepackage{amssymb}
% Para incluir links nas referencias
\usepackage{url}
% Para incluir paginas de documentos .pdf externos
\usepackage{pgfpages}
% Para ajustar o estilo dos contadores
\usepackage{enumerate}
% Para modificar a cor do texto
\usepackage{color}
% Para incluir condicoes
\usepackage{ifthen}
% Para colocar legendas em algo que nao e float
\usepackage{capt-of}

\usepackage{bookmark}
\usepackage[poster]{tcolorbox}
\usepackage[orientation=portrait,size=a4,scale=2.2,debug]{beamerposter}
\usepackage{hyperref}
\usepackage{fontspec}
\setmainfont[UprightFont={* Regular},BoldFont={* Bold},ItalicFont={* Medium}]{AmpleSoft Pro}

\let\sfdefault\rmdefault

% Para definir o tema do slide
\usetheme{Berlin}
\setbeamerfont*{section in head/foot}{size=\large}
\setbeamerfont*{subsection in head/foot}{size=\normalsize}
\setbeamerfont*{author in head/foot}{size=\normalsize}
\setbeamerfont*{title in head/foot}{size=\normalsize}
\setbeamerfont*{date in head/foot}{size=\normalsize}

% Para definir cores e background
\usecolortheme{Asset1}
% Para numerar as figuras
\setbeamertemplate{caption}[numbered]

\def\dt{\VAR{strDate}}
\def\fund{\VAR{strFund}}

\date{\VAR{strLongDate}}

\def\folder{figuras}

\def\traders{
	\VAR{strTraders}}

% Título
\title[Trader Risk Report]{Asset1 Trader Risk Report}

\titlegraphic{
	\includegraphics[width=0.35\textwidth]{logo_Asset1}
}

% Autores
\author[\href{mailto:risco@asset1.com.br}{risco@asset1.com.br}]{
	\href{mailto:riscoa@asset1.com.br}{Team Risk}
	\vspace{0.10cm}
}

\begin{document}	
	
	\foreach \trader\tradername in \traders{
	
	\section[\trader]{\tradername}
	
	\begin{frame}{Risk Breakdown (I) \hspace{10cm} \hbox{\normalsize{\tradername}} }
		\vspace*{5mm}
		\begin{tcbposter}[coverage = {spread}, 
			poster = {%showframe, 
				columns = 1, rows=4, spacing=1mm}, 
			boxes = {enhanced standard jigsaw, sharp corners=northwest,
				arc=1mm,boxrule=0.4mm,colback=white,opacityback=0.75,
				colframe=blue1,	title style={color=blue1}, fonttitle=\bfseries\scshape}
			]
			\posterbox[adjusted title=Value-at-Risk Book Hierarchy, top=0mm, left=0mm, right=0mm, bottom=0mm]
			{name=vargraph, column=1, row=1, rowspan=0.875}
			{ 
				\centering
				\includegraphics[width=0.99\linewidth]{\folder/\dt_\trader_0.pdf}  				
			}	
			\posterbox[adjusted title=Delta Exposure by Book, top=0mm, left=0mm, right=0mm, bottom=0mm]
			{name=deltaexposure, column=1, below=vargraph, rowspan=1.75}
			{ 
			\centering
			\includegraphics[width=0.99\linewidth]{\folder/\dt_\trader_1.pdf} 			
			}	
				
			\posterbox[adjusted title=PnL Historical Simulation, top=0mm, left=0mm, right=0mm, bottom=0mm]
			{name=pnlsimu, column=1, below=deltaexposure, rowspan=1.175}
			{ 
				\centering
			\includegraphics[width=0.99\linewidth]{\folder/\dt_\trader_2.pdf} 
			}
		\end{tcbposter}
	\end{frame}	

	\begin{frame}{Risk Breakdown (II) \hspace{10cm} \hbox{\normalsize{\tradername}} }
		\vspace*{5mm}
		\begin{tcbposter}[coverage = {spread}, 
			poster = {%showframe, 
				columns = 1, rows=3, spacing=1mm}, 
			boxes = {enhanced standard jigsaw, sharp corners=northwest,
				arc=1mm,boxrule=0.4mm,colback=white,opacityback=0.75,
				colframe=blue1,	title style={color=blue1}, fonttitle=\bfseries\scshape}
			]
			\posterbox[adjusted title=Extreme Value Analysis, top=0mm, left=0mm, right=0mm, bottom=0mm]
			{name=evt, column=1, row=1, rowspan=1.2}
			{ 
				\centering
				\includegraphics[width=0.99\linewidth]{\folder/\dt_\trader_evt.pdf}  				
			}	
			
			
			\posterbox[adjusted title={}, top=0mm, left=0mm, right=0mm, bottom=0mm]
			{name=pnlsimu, column=1, below=evt}
			{ 
				\centering
				%\includegraphics[width=0.99\linewidth]{\folder/\dt_\trader_2.pdf} 
			}
		\end{tcbposter}
	\end{frame}	}
	
	

	
\end{document}